% GENERATED FILE. Edit canonical lessons, then run npm run generate:books.
% canonical-source-hash: fnv1a64:cc2c6fc057fce464
% canonical-lessons: ML-C86-oppam, ML-C86-koode, ML-C86-aal-instrumental, ML-R86-with-recall

\chapter{Two Kinds Of With}
\label{ch:ml-with}
\hlchaptermodality{86}

\begin{chapteropening}
\textbf{By the end of this chapter:} \emph{I can say who came along with me and what I did something with, and I know those are different relationships in Malayalam.}

From cold: both company words and their register split, the instrument ending, and the question English never makes a speaker ask.
\end{chapteropening}

\section[oppaṁ]{\ml{ഒപ്പം} (oppaṁ) — together with}

\label{lesson:ML-C86-oppam}



\begin{quote}
Say \emph{on top of the chair}, and say what the chair had to carry first. Then say \emph{my}.
\end{quote}

\subsection*{You'll want to know: \ml{ഒപ്പം}}
\textbf{\ml{ഒപ്പം}} — \emph{oppaṁ} — together with.

You have been able to place a \textbf{thing} since the place-words. This places a \textbf{person}, and it does it in the slot you already know:

\textbf{\ml{എന്റെ} \ml{ഒപ്പം}} — \emph{enṟe oppaṁ} — with me.

Look at what that is. \textbf{\ml{എന്റെ}} is the owner-form you have had since almost the beginning, and \textbf{\ml{ഒപ്പം}} follows it exactly as \textbf{\ml{മുകളിൽ}} followed \emph{the chair's}. Same shape, same order, nothing new to arrange:

\noindent
\begin{tabularx}{\linewidth}{@{}>{\raggedright\arraybackslash}X>{\raggedright\arraybackslash}X@{}}
\toprule
\textbf{} & \textbf{} \\
\midrule
\textbf{\ml{കസേരയുടെ} \ml{മുകളിൽ}} & on top of the chair \\
\textbf{\ml{എന്റെ} \ml{ഒപ്പം}} & together with me \\
\bottomrule
\end{tabularx}

The slot does not care whether what follows it is a \textbf{position} or a \textbf{company}. That is a good deal of reach from one arrangement, and it is the reason the place-words chapter spent itself on the pattern.

Now put a person in it who is not you. \textbf{\ml{സുഹൃത്ത്}} ends in a consonant, so its owner-form takes the ending consonants take:

\textbf{\ml{സുഹൃത്തിന്റെ} \ml{ഒപ്പം}} — \emph{suhṛtthinṟe oppaṁ} — together with a friend.

\subsection*{Your turn}
Say these aloud:

\begin{itemize}
\raggedright
  \item \emph{oppaṁ}
  \item \emph{enṟe oppaṁ}
  \item which earlier word \emph{oppaṁ} is standing in the place of
  \item \emph{suhṛtthinṟe oppaṁ}
\end{itemize}

\begin{tcolorbox}[breakable,colback=teal!4,colframe=teal!35!black,title={Before you move on}]
What does \emph{oppaṁ} mean? (\textbf{Together with}.) What must the word in front of it carry? (\textbf{The owner-ending}.) And where have you seen that arrangement before? (\textbf{With the place-words}.)
\end{tcolorbox}
\section[kūṭe]{\ml{കൂടെ} (kūṭe) — with, alongside}

\label{lesson:ML-C86-koode}



\begin{quote}
Say \emph{with me}, and say what came before the new word in it.
\end{quote}

\subsection*{You'll want to know: \ml{കൂടെ}}
\textbf{\ml{കൂടെ}} — \emph{kūṭe} — with, alongside.

It goes exactly where \emph{oppaṁ} goes and means the same thing:

\textbf{\ml{എന്റെ} \ml{കൂടെ}} — \emph{enṟe kūṭe} — with me.

So why two? \textbf{You have met this arrangement before.} When the degree words arrived, one of them belonged to the page and the other to the room, and the difference was register rather than meaning. This is the same split:

\noindent
\begin{tabularx}{\linewidth}{@{}>{\raggedright\arraybackslash}X>{\raggedright\arraybackslash}X@{}}
\toprule
\textbf{} & \textbf{} \\
\midrule
\textbf{\ml{ഒപ്പം}} \emph{oppaṁ} & the more formal of the two — writing, and careful speech \\
\textbf{\ml{കൂടെ}} \emph{kūṭe} & the everyday one you will hear far more often \\
\bottomrule
\end{tabularx}

Neither is wrong anywhere. \textbf{A learner who owns only one will be understood, and a learner who owns both will place themselves correctly}, which is the same return the degree pair paid.

\subsection*{Your turn}
Say these aloud:

\begin{itemize}
\raggedright
  \item \emph{kūṭe}
  \item \emph{enṟe kūṭe}
  \item \emph{enṟe oppaṁ}, then \emph{enṟe kūṭe}, and say which you would write
  \item the earlier pair that split the same way
\end{itemize}

\begin{tcolorbox}[breakable,colback=teal!4,colframe=teal!35!black,title={Before you move on}]
What does \emph{kūṭe} mean? (\textbf{With, alongside}.) Which of the two is the everyday one? (\emph{\textbf{kūṭe}}.) And does the word in front of it change? (\textbf{No — the owner-ending either way}.)
\end{tcolorbox}
\section[kattiyāl]{\ml{കത്തിയാൽ} (kattiyāl) — by means of}

\label{lesson:ML-C86-aal-instrumental}



\begin{quote}
Say \emph{if he comes}, and say which ending made it \emph{if}. Then say the word for a knife.
\end{quote}

\begin{grammarlens}[title={Grammar Lens: the same ending on a different kind of word}]
\textbf{\ml{ആൽ}} is an ending you already own. On a \textbf{verb} it means \emph{if}.

\textbf{Put it on a noun instead and it means \emph{by means of}.}

\textbf{\ml{കത്തിയാൽ}} — \emph{kattiyāl} — by means of a knife, with a knife.

\noindent
\begin{tabularx}{\linewidth}{@{}>{\raggedright\arraybackslash}X>{\raggedright\arraybackslash}X@{}}
\toprule
\textbf{what it lands on} & \textbf{what it does} \\
\midrule
a \textbf{verb} — \emph{vannāl} & \textbf{if} he comes \\
a \textbf{noun} — \emph{kattiyāl} & \textbf{by means of} a knife \\
\bottomrule
\end{tabularx}

\textbf{The ending has not changed and the meanings are not related.} What decides is the kind of word underneath it, and a reader who knows that will never have to choose between the two readings — the word in front already chose.

This is worth meeting on purpose rather than by surprise. \textbf{One ending doing two unrelated jobs is ordinary in this language}, and the habit it should build is to look at the host before reading the ending.
\end{grammarlens}

\begin{culture}
English hides a distinction here that Malayalam marks. \emph{I went with a friend} and \emph{I cut it with a knife} use the same English word, and \textbf{they are not the same relationship at all}: one is company, the other is a tool.

\noindent
\begin{tabularx}{\linewidth}{@{}>{\raggedright\arraybackslash}X>{\raggedright\arraybackslash}X@{}}
\toprule
\textbf{} & \textbf{} \\
\midrule
\textbf{company} — a person who came along & \textbf{\ml{ഒപ്പം}}, after the owner-form \\
\textbf{instrument} — a thing you used & \textbf{-\ml{ആൽ}}, glued onto the noun \\
\bottomrule
\end{tabularx}

So the question to ask is never \emph{how do I say \textquotedblleft{}with\textquotedblright{}}. It is \textbf{which kind of \emph{with} this is} — and once you ask that, the two never compete.

\textbf{-\ml{ആൽ} belongs more to writing than to talk.} You will read it more often than you say it, which is a good reason to recognise it and no reason to avoid it.
\end{culture}

\subsection*{Your turn}
Say these aloud:

\begin{itemize}
\raggedright
  \item \emph{kattiyāl}
  \item \emph{vannāl}, then \emph{kattiyāl}, and say what kind of word each ending sat on
  \item which of the two \textquotedblleft{}with\textquotedblright{}s a friend takes, and which a knife takes
  \item which of the two you are more likely to read than to say
\end{itemize}

\begin{tcolorbox}[breakable,colback=teal!4,colframe=teal!35!black,title={Before you move on}]
What does \textbf{-\ml{ആൽ}} mean on a noun? (\textbf{By means of}.) What does it mean on a verb? (\emph{\textbf{If}}.) And what tells you which reading to take? (\textbf{The kind of word it is sitting on}.)
\end{tcolorbox}
\section[kūṭe enna ōrmma]{(two kinds of with) — from cold}

\label{lesson:ML-R86-with-recall}



\begin{quote}
Close the lessons before this one. Say \emph{with me} twice, once with each of the two words.
\end{quote}

\begin{grammarlens}[title={Grammar Lens: English asks one question, Malayalam asks two}]
\emph{I went \textbf{with} a friend.} \emph{I cut it \textbf{with} a knife.} \textbf{One English word, two relationships that have nothing in common} — a companion and a tool.

\noindent
\begin{tabularx}{\linewidth}{@{}>{\raggedright\arraybackslash}X>{\raggedright\arraybackslash}X>{\raggedright\arraybackslash}X@{}}
\toprule
\textbf{} & \textbf{} & \textbf{} \\
\midrule
\textbf{company} & \textbf{\ml{എന്റെ} \ml{ഒപ്പം}} / \textbf{\ml{എന്റെ} \ml{കൂടെ}} & after the owner-form, standing apart \\
\textbf{instrument} & \textbf{\ml{കത്തിയാൽ}} & glued onto the noun \\
\bottomrule
\end{tabularx}

So the useful question is never \emph{how do I say with}. It is \textbf{which kind of \emph{with} this is.} Ask that and the two never compete; skip it and you will pick by coin-toss.

The company pair splits by register the way the degree words did — \textbf{\ml{ഒപ്പം}} to the page, \textbf{\ml{കൂടെ}} to the room — and both sit in the slot the place-words established, after the owner-ending and unchanged.
\end{grammarlens}

\begin{grammarlens}[title={What you've built}]
A distinction your own language does not draw for you.

And a second look at an ending you thought you had finished with. \textbf{\ml{ആൽ}} on a \textbf{verb} is \emph{if}; on a \textbf{noun} it is \emph{by means of}. The ending is identical and the meanings are unrelated — \textbf{what decides is the kind of word underneath it.}

That is the habit worth taking away: \textbf{look at the host before you read the ending.} It will keep paying, because one ending doing two unrelated jobs is ordinary here rather than exceptional.
\end{grammarlens}

\subsection*{Your turn}
Say these aloud:

\begin{itemize}
\raggedright
  \item \emph{enṟe oppaṁ}, then \emph{enṟe kūṭe}
  \item \emph{suhṛtthinṟe kūṭe}
  \item \emph{kattiyāl}
  \item which of the three is a tool and which are companions
  \item what decides whether \textbf{\ml{ആൽ}} means \emph{if} or \emph{by}
\end{itemize}

\begin{tcolorbox}[breakable,colback=teal!4,colframe=teal!35!black,title={Before you move on}]
Which word do you use for a person who came along? (\emph{\textbf{oppaṁ}} or \emph{\textbf{kūṭe}}.) Which ending marks a tool? (\textbf{-\ml{ആൽ}}.) And what question should you ask before translating \emph{with}? (\textbf{Which kind of \emph{with} is this}.)
\end{tcolorbox}
